\documentclass{article}
\usepackage[utf8]{inputenc}
\usepackage{amsmath}
\usepackage{amsfonts}
\usepackage{array}
\usepackage{booktabs}
\usepackage[spanish]{babel}
\usepackage[left=1cm, right=1cm, top=1.5cm, bottom=1.5cm]{geometry}

\title{Formulario de Fórmulas de Estadística}
\date{}

\begin{document}

\maketitle

\section{Teoría de Conjuntos}
Estos son los conceptos básicos para manejar colecciones de elementos.

\subsection{Pertenencia y Subconjunto}
\( x \in A \): Indica que \( x \) es un elemento de \( A \). \\
\( A \subseteq B \): \( A \) es un subconjunto de \( B \) (todos los elementos de \( A \) están en \( B \)). \\
\textit{Explicación}: Define relaciones básicas entre elementos y conjuntos. El conjunto vacío \( \emptyset \subseteq A \) para cualquier \( A \).

\subsection{Cardinal de un Conjunto}
\( |A| \) o \( \#(A) \): Número de elementos en \( A \). \\
\textit{Explicación}: Mide el tamaño finito de un conjunto. Ejemplo: \( |\{1, 2, 3\}| = 3 \).

\subsection{Conjunto Potencia (Partes)}
\( \mathcal{P}(A) \): Conjunto de todos los subconjuntos de \( A \). \\
\( |\mathcal{P}(A)| = 2^{|A|} \): Cardinal del conjunto potencia. \\
\textit{Explicación}: Incluye todos los subconjuntos posibles, incluyendo \( \emptyset \) y \( A \). Ejemplo: Para \( A = \{1, 2\} \), \( |\mathcal{P}(A)| = 4 \).

\subsection{Función Característica}
\[
\chi_A(\omega) = 
\begin{cases} 
1 & \text{si } \omega \in A \\ 
0 & \text{si } \omega \notin A 
\end{cases}
\]
\textit{Explicación}: Asigna 1 si el elemento está en \( A \), 0 en caso contrario. Útil para operaciones booleanas.

\subsection{Operaciones Básicas}
\begin{itemize}
    \item Intersección: \( A \cap B = \{ x \in \Omega \mid x \in A \ y \ x \in B \} \). (Elementos comunes).
    \item Unión: \( A \cup B = \{ x \in \Omega \mid x \in A \ o \ x \in B \} \). (Elementos en al menos uno).
    \item Diferencia: \( A - B = \{ x \in \Omega \mid x \in A \ y \ x \notin B \} \). (Elementos en \( A \) pero no en \( B \)).
    \item Complemento: \( A^c = \Omega - A = \{ x \in \Omega \mid x \notin A \} \). (Elementos no en \( A \)).
\end{itemize}
\textit{Explicación}: Operaciones fundamentales para manipular conjuntos. Propiedades: conmutativas (\( A \cup B = B \cup A \)), asociativas, distributivas.

\subsection{Leyes de De Morgan}
\( {(A \cup B)}^c = A^c \cap B^c \). \\
\( {(A \cap B)}^c = A^c \cup B^c \). \\
\textit{Explicación}: Relacionan complementos con uniones e intersecciones. Útil para simplificar expresiones.

\subsection{Propiedades Adicionales}
Dos conjuntos son disjuntos si \( A \cap B = \emptyset \). \\
\( \emptyset^c = \Omega \), \( \Omega^c = \emptyset \). \\
\textit{Explicación}: Reglas básicas para conjuntos vacíos y universales.

\section{Combinatoria}
Se enfoca en contar configuraciones de elementos.

\subsection{Número Binomial (Combinaciones sin Repetición)}
\[
\binom{n}{k} = C_k^n = \frac{n!}{k! \cdot (n-k)!}
\]
\textit{Explicación}: Número de subconjuntos de tamaño \( k \) en un conjunto de \( n \) elementos (orden no importa, sin repetición). Ejemplo: \( \binom{7}{5} = 21 \) (equipos de 5 con 7 posibles jugadores).

\subsection{Combinaciones con Repetición}
\[
CR_k^n = \binom{n + k - 1}{k} = \frac{(n + k - 1)!}{k! \cdot (n-1)!}
\]
\textit{Explicación}: Número de formas de elegir \( k \) elementos de \( n \) tipos permitiendo repeticiones (multiconjuntos). Ejemplo: Repartir 12 caramelos en 4 personas: \( \binom{4 + 12 - 1}{12} = 455 \).

\subsection{Variaciones sin Repetición}
\[
V_k^n = \frac{n!}{(n-k)!} = (n - k + 1) \cdot (n - k + 2) \cdots n
\]
\textit{Explicación}: Número de formas de elegir y ordenar \( k \) elementos de \( n \) sin repetir (permutaciones parciales). Ejemplo: Números de 2 cifras de \{1,2,3\}: \( V_2^3 = 6 \).

\subsection{Variaciones con Repetición}
\( VR_k^n = n^k \). \\
\textit{Explicación}: Similar a variaciones, pero permitiendo repeticiones. Ejemplo: Números de 2 cifras de \{1,2,3\} con repetición: \( 3^2 = 9 \).

\subsection{Permutaciones (sin Repetición)}
\( P_n = n! \). \\
\textit{Explicación}: Número de formas de ordenar todos los \( n \) elementos de un conjunto. Ejemplo: Permutaciones de \{1,2,3\}: \( 3! = 6 \).

\subsection{Permutaciones con Repetición (Multinomial)}
\[
PR_{n_1, n_2, \dots, n_k}^n = \binom{n}{n_1 n_2 \dots n_k} = \frac{n!}{n_1! \cdot n_2! \cdots n_k!}, \quad \text{donde } n = n_1 + \cdots + n_k
\]
\textit{Explicación}: Número de permutaciones cuando hay elementos repetidos. Ejemplo: Palabras con letras de ``PROBABILIDAD'' (12 letras, repeticiones específicas): \( \frac{12!}{{2!}^4 {1!}^4} = 29,937,600 \).

\subsection{Principio de la Suma}
Si \( A_1, \dots, A_n \) son disjuntos: \( \#(\cup_{i=1}^n A_i) = \sum_{i=1}^n \#(A_i) \). \\
\textit{Explicación}: Suma los cardinales de conjuntos mutuamente exclusivos.

\subsection{Principio de Inclusión-Exclusión (para 2 Conjuntos)}
\( \#(A_1 \cup A_2) = \#(A_1) + \#(A_2) - \#(A_1 \cap A_2) \). \\
\textit{Explicación}: Corrige el solapamiento al contar uniones de conjuntos no disjuntos.

\subsection{Principio del Producto}
\( \#(A_1 \times A_2 \times \cdots \times A_n) = \prod_{i=1}^n \#(A_i) \). \\
\textit{Explicación}: Cardinal de un producto cartesiano (selecciones independientes).

\section{Conceptos Básicos y Operaciones con Sucesos}

\subsection{Definiciones Básicas}
\begin{itemize}
    \item \textbf{Experimento aleatorio}: Experimento que, repetido en las mismas condiciones, puede dar resultados diferentes, pero predecibles a largo plazo.
    \item \textbf{Suceso elemental}: Cada posible resultado del experimento.
    \item \textbf{Espacio muestral} \(\Omega\): Conjunto de todos los sucesos elementales.
    \item \textbf{Suceso}: Cualquier subconjunto de \(\Omega\).
    \item \textbf{Suceso seguro}: \(\Omega\), \(P(\Omega) = 1\).
    \item \textbf{Suceso imposible}: \(\emptyset\), \(P(\emptyset) = 0\).
    \item \textbf{Conjunto de partes} \(\mathcal{P}(\Omega)\): Todos los subconjuntos de \(\Omega\). Si \(|\Omega| = n\), entonces \(|\mathcal{P}(\Omega)| = 2^n\).
\end{itemize}

\subsection{Operaciones con Sucesos}
Sea \(A, B \subseteq \Omega\):
\begin{itemize}
    \item \textbf{Unión}: \(A \cup B\) (ocurre si sucede \(A\) o \(B\)).
    \item \textbf{Intersección}: \(A \cap B\) (ocurre si sucede \(A\) y \(B\)).
    \item \textbf{Complementario}: \(A^c\) (ocurre si no sucede \(A\)).
    \item \textbf{Diferencia}: \(A - B = A \cap B^c\) (ocurre si sucede \(A\) y no \(B\)).
    \item \textbf{Sucesos incompatibles (disjuntos)}: \(A \cap B = \emptyset\).
\end{itemize}

\subsection{Propiedades de Operaciones con Sucesos}
\begin{itemize}
    \item \textbf{Conmutativas}: \(A \cup B = B \cup A\), \(A \cap B = B \cap A\).
    \item \textbf{Asociativas}: \(A \cup (B \cup C) = (A \cup B) \cup C\), \(A \cap (B \cap C) = (A \cap B) \cap C\).
    \item \textbf{Distributivas}: \(A \cap (B \cup C) = (A \cap B) \cup (A \cap C)\), \(A \cup (B \cap C) = (A \cup B) \cap (A \cup C)\).
    \item \textbf{Complementario del complementario}: \({(A^c)}^c = A\).
    \item \textbf{Leyes de De Morgan}: \({(A \cup B)}^c = A^c \cap B^c\), \({(A \cap B)}^c = A^c \cup B^c\).
\end{itemize}

\subsection{Definición de Probabilidad}
Para un espacio muestral finito \(\Omega\), una probabilidad \(P: \mathcal{P}(\Omega) \to [0,1]\) satisface:
\begin{enumerate}
    \item \(0 \leq P(A) \leq 1\) para todo suceso \(A\).
    \item \(P(\Omega) = 1\).
    \item Si \(A_1, \dots, A_n\) disjuntos, \(P(A_1 \cup \cdots \cup A_n) = P(A_1) + \cdots + P(A_n)\).
\end{enumerate}
Si todos los elementales son equiprobables, \(P(A) = \frac{|A|}{|\Omega|}\).

\section{Probabilidad}

\subsection{Propiedades}
\begin{itemize}
    \item \(P(\emptyset) = 0\).
    \item \(P(A^c) = 1 - P(A)\).
    \item Si \(B \subseteq A\), entonces \(0 \leq P(B) \leq P(A)\).
    \item \(P(A - B) = P(A) - P(A \cap B)\).
    \item \(P(A \cup B) = P(A) + P(B) - P(A \cap B)\).
    \item Para tres conjuntos: \(P(A \cup B \cup C) = P(A) + P(B) + P(C) - P(A \cap B) - P(A \cap C) - P(B \cap C) + P(A \cap B \cap C)\).
\end{itemize}

\subsection{Probabilidad Condicionada}
Dado \(A, B\) con \(P(A) > 0\):
\[
P(B \mid A) = \frac{P(A \cap B)}{P(A)}.
\]
Propiedades (similar a probabilidad estándar):
\begin{itemize}
    \item \(P(B^c \mid A) = 1 - P(B \mid A)\).
    \item \(P(B_1 \cup B_2 \mid A) = P(B_1 \mid A) + P(B_2 \mid A) - P(B_1 \cap B_2 \mid A)\).
\end{itemize}

\subsection{Teorema de la Probabilidad Total}
Para \(A, B\):
\[
P(B) = P(B \cap A) + P(B \cap A^c) = P(A) \cdot P(B \mid A) + P(A^c) \cdot P(B \mid A^c).
\]
Para partición \(\{A_1, \dots, A_n\}\) de \(\Omega\):
\[
P(B) = \sum_{i=1}^n P(A_i) \cdot P(B \mid A_i).
\]

\subsection{Teorema de Bayes}
Para \(A, B\) con \(P(B) > 0\):
\[
P(A \mid B) = \frac{P(A) \cdot P(B \mid A)}{P(B)} = \frac{P(A) \cdot P(B \mid A)}{P(A) \cdot P(B \mid A) + P(A^c) \cdot P(B \mid A^c)}.
\]
Para partición \(\{A_1, \dots, A_n\}\):
\[
P(A_i \mid B) = \frac{P(A_i) \cdot P(B \mid A_i)}{\sum_{j=1}^n P(A_j) \cdot P(B \mid A_j)}.
\]

\subsection{Independencia de Sucesos}
\(A\) y \(B\) son independientes si:
\[
P(A \cap B) = P(A) \cdot P(B).
\]
Equivalentes (si \(P(A), P(B) > 0\)):
\begin{itemize}
    \item \(P(A \mid B) = P(A)\).
    \item \(P(B \mid A) = P(B)\).
\end{itemize}
Para familia \(\{A_1, \dots, A_n\}\): Para toda subfamilia, el producto de probabilidades individuales iguala la probabilidad de la intersección.

También: Independencia se preserva con complementarios (e.g., \(A^c\) y \(B\) independientes si \(A\) y \(B\) lo son).

\end{document}